%%%%%
%
\section{Fonction exponentielle polynomiale dans un domaine de d\'efinition}

\itemstyle{Op\'erateur~:} \textsf{Poly\_exp\_domain}
\index{Poly\_exp\_domain}

%%%%%OBJET%%%%
\itemstyle{Objet~:}
Application d'une fonction exponentielle polynomiale (de degr\'e 2) appliqu\'ee dans un domaine de d\'efinition. Si la variable en entr\'ee n'appartient pas au domaine de d\'efinition (valeur comprise entre borne mini et borne maxi), alors l'op\'erateur renvoie une valeur manquante.

%%%%%ARGUMENTS%%%%
\bigskip
%%%%%ENTREES%%%%
\begin{tabarg}
\multicolumn{3}{l}{\titretab {Entr\'ees}}\\
\hline
$X$ & Vecteur, {\tt float} & Op\'erande (unit\'e quelconque)\\
$C_0$ & Coeff., {\tt double} & Coefficient du terme de degr\'e 0.\\
$C_1$ & Coeff., {\tt double} & Coefficient du terme de degr\'e 1.\\
$C_2$ & Coeff., {\tt double} & Coefficient pour l'exposant du terme de degr\'e 2.\\
$C_3$ & Coeff., {\tt double} & Coefficient pour l'exposant du terme de degr\'e 1.\\
$C_4$ & Coeff., {\tt double} & Coefficient pour l'exposant du terme de degr\'e 0.\\
$Borne\_{min}$ & Coeff., {\tt float} & Borne minimale pour le parametre X\\
$Borne\_{max}$ & Coeff., {\tt float} & Borne maximale pour le parametre X\\
$Unit\_{res}$ & Coeff., {\tt String} & Unit\'e du r\'esultat.\\
\end{tabarg}

\medskip
%%%%%SORTIES%%%%
\begin{tabarg}
\multicolumn{3}{l}{\titretab {Sorties}}\\
\hline
$Z$ & Vecteur, {\tt float} & R\'esultat (unit\'e indiqu\'ee par $Unit\_{res}$)\\
\end{tabarg}


%%%%%FORMULE%%%%
\itemstyle{Formule}

%% Poly_exp_domain
Uniquement si la variable X est dans le domaine (sup\'erieure ou \'egale \`a Borne\_mini et inf\'erieure ou \'egale \`a Borne\_maxi) : \\
On calcule d'abord une variable interm\'ediaire Y qui est une fonction polynomiale de X (degr\'e 2) : \\
\hspace*{2cm} $Y = C_2 + C_3 . X + C_4 . X^2$ \\
Puis, on calcule le r\'esultat Z , fonction exponentielle de $Y$ : \\
\hspace*{2cm} $Z = C_0 + C_1 . e^Y$ \\



%%%%%COMMENTAIRE%%%%
\itemstyle{Commentaire~:}
On effectue le calcul uniquement si la variable X est dans le domaine de d\'efinition :
 c'est \`a dire sup\'erieure ou \'egale \`a Borne\_mini et inf\'erieure ou \'egale \`a Borne\_maxi.

%%%%%AUTEUR%%%%
\itemstyle{Auteur~:}
CNRM/GMEI/TRAMM

%%%%%REFERENCE%%%%
\itemstyle{R\'ef\'erences~:}
